\begin{abstract}
A false answer can reflect uncertainty, a stable misconception, or misreporting under an incentive. A useful deception detector must distinguish these possibilities without relying only on the surrounding prompt. We conduct a matched-question diagnostic study of a fixed historical model, Qwen2.5-7B-Instruct, using 600 unique TriviaQA questions and 11,400 local generations. Independent knowledge checks, five response conditions, and context-preserving follow-ups define conservative operational labels. Only two incentive-associated false responses satisfy our recovered-misreport criterion; one is a malformed fragment. Across an equal mixture of four non-instructed conditions, these proxies account for 0.21\% of 939 strictly scored false outputs, while 60.49\% remain unresolved. A reference-equivalence audit yields 0.25\% and 62.36\%, respectively. An instructed-deception response probe achieves AUROC 0.972 against uncertain neutral errors, but a pre-answer prompt-only probe achieves 1.000. The response probe reaches only 0.605 on unverified incentive mismatches and flags 22.6\% of known-correct reward answers at a neutral-calibrated threshold, compared with 3.7\% of known-correct neutral answers. With only one recovered test example, the central verified-misreport comparison remains underpowered. Our results identify prompt sensitivity and calibration failures in this protocol, while showing why high instructed-condition discrimination does not establish detection of belief--statement mismatch.
\end{abstract}
